\chapter{Методы градиента стратегии}
\label{ch:pg}

\begin{quote}
\itshape
Вместо того чтобы искать стратегию косвенно через функцию ценности, мы ищем её напрямую. Мы спрашиваем: при данной текущей стратегии, в каком направлении следует скорректировать параметры, чтобы ожидаемое вознаграждение возросло? Этот вопрос имеет прекрасный ответ, известный под названием теоремы о градиенте стратегии.
\end{quote}

\bigskip

Все методы, изученные в главах~\ref{ch:dp}--\ref{ch:value-approx}, разделяют общую архитектуру: сначала оценить функцию ценности, затем вывести из неё стратегию, действуя жадно или $\varepsilon$-жадно. Функции ценности являются полезными промежуточными объектами, но они не строго необходимы. Вместо этого можно параметризовать стратегию напрямую\,---\,как дифференцируемую функцию $\pi_\theta(a \mid s)$ от параметров $\theta$\,---\,и оптимизировать эти параметры для максимизации ожидаемой отдачи. Это идея \textbf{градиента стратегии}\index{policy gradient}, и именно она является предметом данной главы.

Методы градиента стратегии имеют три важных преимущества перед методами, основанными на функции ценности. Во-первых, они естественно обрабатывают непрерывные и многомерные пространства действий: выборка из $\pi_\theta(\cdot \mid s)$ не требует внутренней максимизации. Во-вторых, обученная стратегия может быть стохастической, что иногда является доказуемо оптимальным (например, в частично наблюдаемых средах или в играх с противником) и всегда является корректно определённым вероятностным распределением. В-третьих, и пожалуй наиболее важным для современных приложений, градиент ожидаемой отдачи по параметрам стратегии может быть оценён исключительно по выборкам траекторий, без необходимости дифференцировать через динамику среды.

Глава организована следующим образом. Раздел~\ref{sec:pg-why} объясняет мотивацию прямой оптимизации стратегии. Раздел~\ref{sec:pg-theorem} выводит теорему о градиенте стратегии\,---\,математическую основу всего семейства методов\,---\,включая полные доказательства трюка логарифмического отношения, сведения к отдаче с текущего момента с помощью принципа причинности, и компактной формы через $Q^{\pi_\theta}$. Раздел~\ref{sec:reinforce} представляет алгоритм REINFORCE с разобранным числовым примером. Раздел~\ref{sec:baselines} доказывает, что вычитание базовой линии, зависящей от состояния, не вносит смещения в градиент, и выводит оптимальную базовую линию. Раздел~\ref{sec:reinforce-baseline} объединяет оценку функции ценности с градиентом стратегии, получая REINFORCE с обученной базовой линией. Раздел~\ref{sec:rloo} разрабатывает REINFORCE Leave-One-Out (RLOO)\,---\,современный оценщик, играющий ключевую роль в выравнивании больших языковых моделей. Раздел~\ref{sec:npg} вводит натуральный градиент стратегии, который корректирует обычный градиент с учётом кривизны многообразия стратегий. Раздел~\ref{sec:pg-convergence} анализирует сходимость и компромисс между смещением и дисперсией. Раздел~\ref{sec:pg-comparison} сравнивает оценщики в итоговой таблице. Глава завершается упражнениями.

% ===========================================================
\section{Зачем оптимизировать стратегию напрямую?}
\label{sec:pg-why}

\subsection{Ограничения подхода, основанного на функции ценности}

Пайплайн подхода, основанного на функции ценности: обучить $\hat Q \approx Q^*$, затем задать $\pi(s) = \argmax_a \hat Q(s,a)$. Этот пайплайн ломается двумя важными способами, когда пространство действий велико или непрерывно.

\begin{enumerate}[leftmargin=*]
\item \textit{Непрерывные действия.}  Если $\cA \subseteq \R^m$, вычисление $\argmax_{a \in \R^m} \hat Q(s,a)$ само по себе является задачей непрерывной оптимизации, которую необходимо решать на каждом шаге во время выполнения. Такие методы, как DDPG, обходят эту проблему, обучая отдельную сеть-максимизатор, но в результате получается архитектура актор-критик\,---\,которая уже выходит за рамки чистого обучения на основе функции ценности.

\item \textit{Очень большие дискретные словари.}  В генерации языка на уровне токенов $|\cA|$ равен размеру словаря (часто $50{,}000$ или более), а длина эпизода может составлять сотни шагов. Хранить отдельное значение для каждой пары $(s,a)$ невозможно; и даже при использовании аппроксимации функции, перебор всех действий для нахождения жадного является дорогостоящим.
\end{enumerate}

Обе трудности исчезают, когда стратегия является параметризованным распределением $\pi_\theta(a \mid s)$: во время выполнения достаточно просто выбрать $A_t \sim \pi_\theta(\cdot \mid S_t)$, независимо от размера $\cA$.

\subsection{Стохастические стратегии иногда необходимы}

Помимо масштабируемости, есть принципиальные причины желать стохастических стратегий. В частично наблюдаемых средах оптимальная стратегия, как правило, стохастична: агенту может потребоваться рандомизация для хеджирования от скрытых состояний. В играх с противником детерминированная стратегия поддаётся эксплуатации, тогда как равновесие смешанных стратегий требует рандомизации. При исследовании стохастическая стратегия выбирает действия без необходимости отдельного механизма $\varepsilon$-жадности, а степень исследования автоматически контролируется энтропией $\pi_\theta$.

\begin{remark}
Детерминированная стратегия $\mu_\theta(s)$ является частным случаем: она соответствует распределению, сосредоточенному в одном действии. Теорема о \emph{детерминированном градиенте стратегии} (DPG) Сильвера и~др.\  (2014) расширяет разработанный здесь аппарат на этот вырожденный случай и приводит к таким алгоритмам, как DDPG и TD3. Мы не будем развивать это направление дальше; оно рассматривается в главе~\ref{ch:actor-critic}.
\end{remark}

\subsection{Целевая функция}

\index{policy gradient!objective}
В эпизодической задаче с горизонтом $T$ мы оптимизируем ожидаемую дисконтированную отдачу:
\begin{equation}
\label{eq:pg-objective}
J(\theta) = \E_{\tau \sim p_\theta}\!\left[\sum_{t=0}^{T-1} \gamma^t R_{t+1}\right],
\end{equation}
где траектория $\tau = (S_0, A_0, R_1, S_1, A_1, R_2, \ldots, S_{T-1}, A_{T-1}, R_T)$ распределена согласно
\begin{equation}
\label{eq:trajectory-density}
p_\theta(\tau) = \rho_0(S_0)\prod_{t=0}^{T-1} \pi_\theta(A_t \mid S_t)\, p(S_{t+1}, R_{t+1} \mid S_t, A_t),
\end{equation}
$\rho_0$ \,---\, это распределение начального состояния, а $p$ \,---\, (неизвестное) ядро переходов. В задаче без эпизодов целевая функция является средним вознаграждением $J(\theta) = \lim_{T\to\infty} \frac{1}{T}\E[\sum_{t=0}^{T-1} R_{t+1}]$; мы сосредотачиваемся на эпизодическом случае для ясности, отмечая, что теорема о градиенте стратегии справедлива в обоих случаях.

Центральный вопрос: как меняется $J(\theta)$ при изменении $\theta$? Наивное дифференцирование~\eqref{eq:pg-objective} кажется требующим дифференцирования через неизвестную динамику среды $p$. Теорема о градиенте стратегии показывает, что это не нужно.

% ===========================================================
\section{Теорема о градиенте стратегии}
\label{sec:pg-theorem}

\index{policy gradient theorem}

\subsection{Трюк логарифмического отношения}

Начнём с вычисления $\nabla_\theta J(\theta)$ непосредственно из определения. Поскольку горизонт конечен, можно дифференцировать под знаком математического ожидания:
\begin{align}
\nabla_\theta J(\theta)
&= \nabla_\theta \sum_\tau p_\theta(\tau)\, R(\tau)
= \sum_\tau \nabla_\theta p_\theta(\tau)\, R(\tau),
\label{eq:pg-raw}
\end{align}
где $R(\tau) = \sum_{t=0}^{T-1} \gamma^t R_{t+1}$ \,---\, полная дисконтированная отдача траектории $\tau$. Эта сумма не является математическим ожиданием, поэтому её нельзя оценить по выборкам\,---\,пока что. \textbf{Трюк логарифмического отношения}\index{likelihood-ratio trick} (также называемый \emph{трюком функции счёта} или \emph{трюком REINFORCE}) преобразует его в таковое. Для любого $p_\theta(\tau) > 0$:
\begin{equation}
\label{eq:log-derivative}
\nabla_\theta p_\theta(\tau) = p_\theta(\tau)\, \nabla_\theta \log p_\theta(\tau).
\end{equation}
Подставляя в~\eqref{eq:pg-raw}, получаем
\begin{equation}
\label{eq:pg-score}
\nabla_\theta J(\theta) = \E_{\tau \sim p_\theta}\!\bigl[R(\tau)\, \nabla_\theta \log p_\theta(\tau)\bigr].
\end{equation}

Теперь воспользуемся факторизованной формой~\eqref{eq:trajectory-density}. Логарифмируя:
\[
\log p_\theta(\tau) = \log \rho_0(S_0) + \sum_{t=0}^{T-1} \log \pi_\theta(A_t \mid S_t) + \sum_{t=0}^{T-1} \log p(S_{t+1}, R_{t+1} \mid S_t, A_t).
\]
Первый и третий члены не зависят от $\theta$, поэтому их градиенты обращаются в нуль:
\begin{equation}
\label{eq:log-traj-grad}
\nabla_\theta \log p_\theta(\tau) = \sum_{t=0}^{T-1} \nabla_\theta \log \pi_\theta(A_t \mid S_t).
\end{equation}
Принципиально важно, что эта формула включает только \emph{стратегию}, а не ядро переходов. Подставляя в~\eqref{eq:pg-score}, получаем базовый оценщик градиента стратегии:

\begin{theorem}[Базовый оценщик градиента стратегии]
\label{thm:basic-pg}
\index{policy gradient!basic estimator}
При приведённых выше предположениях,
\begin{equation}
\label{eq:basic-pg}
\nabla_\theta J(\theta)
= \E_{\tau \sim p_\theta}\!\left[
  \left(\sum_{t=0}^{T-1} \gamma^t R_{t+1}\right)
  \sum_{t=0}^{T-1} \nabla_\theta \log \pi_\theta(A_t \mid S_t)
\right].
\end{equation}
\end{theorem}

Уравнение~\eqref{eq:basic-pg} уже является корректным оценщиком градиента: для данной одиночной траектории $\tau$ выражение внутри математического ожидания является несмещённой выборкой $\nabla_\theta J(\theta)$. Однако оно имеет серьёзную проблему с дисперсией: полная отдача $R(\tau)$ умножает градиент логарифма вероятности каждого действия в траектории, даже действий, которые произошли \emph{до} получения вознаграждений. Интуитивно, то, что произошло при $t=3$, не может задним числом изменить вознаграждение, полученное при $t=1$. Следующая лемма формализует это.

\subsection{Лемма о причинности и отдача с текущего момента}

\begin{lemma}[Причинность]
\label{lem:causality}
\index{causality lemma}
Для каждого $t \in \{0, \ldots, T-1\}$,
\begin{equation}
\label{eq:causality}
\E_{\tau \sim p_\theta}\!\left[\nabla_\theta \log \pi_\theta(A_t \mid S_t)\sum_{k=0}^{t-1} \gamma^k R_{k+1}\right] = 0.
\end{equation}
\end{lemma}

\begin{proof}
Пусть $C_t = \sum_{k=0}^{t-1} \gamma^k R_{k+1}$ обозначает накопленное дисконтированное вознаграждение \emph{до} момента $t$. Эта величина измерима относительно истории $(S_0, A_0, R_1, \ldots, S_t)$ и, в частности, не зависит от выбора $A_t$ при данном $S_t$. По теореме о полном математическом ожидании,
\[
\E\!\left[\nabla_\theta \log \pi_\theta(A_t \mid S_t)\, C_t\right]
= \E\!\left[C_t\, \E\!\left[\nabla_\theta \log \pi_\theta(A_t \mid S_t) \mid S_t, \mbox{история}_{<t}\right]\right].
\]
Для внутреннего математического ожидания зафиксируем $S_t = s$. Тогда $A_t \sim \pi_\theta(\cdot \mid s)$, и
\begin{align*}
&\E\!\left[\nabla_\theta \log \pi_\theta(A_t \mid S_t) \mid S_t = s\right]
= \sum_a \pi_\theta(a \mid s)\, \nabla_\theta \log \pi_\theta(a \mid s)\\
&\quad= \sum_a \nabla_\theta \pi_\theta(a \mid s)
= \nabla_\theta \!\sum_a \pi_\theta(a \mid s)
= \nabla_\theta 1 = 0.
\end{align*}
Поскольку внутреннее условное математическое ожидание тождественно равно нулю, внешнее ожидание также равно нулю.
\end{proof}

\begin{remark}
Ключевое вычисление $\sum_a \pi_\theta(a \mid s)\, \nabla_\theta \log \pi_\theta(a \mid s) = \nabla_\theta \sum_a \pi_\theta(a \mid s) = 0$ называется \textbf{тождеством функции счёта}\index{score function identity}. Это алгебраический механизм, лежащий в основе как теоремы о базовой линии, так и леммы о причинности.
\end{remark}

Определим \textbf{отдачу с текущего момента}\index{reward-to-go} начиная с шага $t$ как
\begin{equation}
\label{eq:rtg}
G_t = \sum_{k=t}^{T-1} \gamma^{k-t} R_{k+1}, \qquad \bar G_t = \gamma^t G_t = \sum_{k=t}^{T-1} \gamma^k R_{k+1}.
\end{equation}
Полная отдача разлагается как $R(\tau) = C_t + \bar G_t$. Применяя лемму~\ref{lem:causality} к~\eqref{eq:basic-pg}, получаем редуцированный оценщик:

\begin{theorem}[Градиент стратегии через отдачу с текущего момента]
\label{thm:rtg-pg}
\index{policy gradient!reward-to-go form}
\begin{equation}
\label{eq:rtg-pg}
\nabla_\theta J(\theta)
= \E_{\tau \sim p_\theta}\!\left[\sum_{t=0}^{T-1} \nabla_\theta \log \pi_\theta(A_t \mid S_t)\, \bar G_t\right]
= \E_{\tau \sim p_\theta}\!\left[\sum_{t=0}^{T-1} \gamma^t\, \nabla_\theta \log \pi_\theta(A_t \mid S_t)\, G_t\right].
\end{equation}
\end{theorem}

\begin{proof}
Начиная с теоремы~\ref{thm:basic-pg}, меняем порядок суммирования по $t$ и суммы по вознаграждениям. Для каждого фиксированного $t$ разложим $R(\tau) = C_t + \bar G_t$ и применим линейность математического ожидания:
\[
\E\!\left[\nabla_\theta \log \pi_\theta(A_t \mid S_t)\, R(\tau)\right]
= \underbrace{\E\!\left[\nabla_\theta \log \pi_\theta(A_t \mid S_t)\, C_t\right]}_{=\,0 \text{ по лемме~\ref{lem:causality}}}
+ \E\!\left[\nabla_\theta \log \pi_\theta(A_t \mid S_t)\, \bar G_t\right].
\]
Суммируя по $t$ и используя $\bar G_t = \gamma^t G_t$, получаем~\eqref{eq:rtg-pg}.
\end{proof}

\subsection{Теорема о градиенте стратегии: компактная форма}

Форма с отдачей с текущего момента~\eqref{eq:rtg-pg} является наиболее практичным оценщиком. Её также можно записать в компактной форме по парам состояние-действие через усреднение по дисконтированной мере посещаемости состояний.

\begin{definition}[Дисконтированная мера посещаемости]
\label{def:visitation}
\index{discounted visitation measure}
\[
d^{\gamma}_{\pi_\theta}(s) = \sum_{t=0}^{T-1} \gamma^t \Prob_{\pi_\theta}(S_t = s).
\]
\end{definition}

\begin{theorem}[Теорема о градиенте стратегии]
\label{thm:pgt}
\index{policy gradient theorem}
\begin{equation}
\label{eq:pgt}
\nabla_\theta J(\theta)
= \sum_{s \in \cS} d^{\gamma}_{\pi_\theta}(s) \sum_{a \in \cA} \nabla_\theta \pi_\theta(a \mid s)\, Q^{\pi_\theta}(s,a)
= \E_{S \sim d^{\gamma}_{\pi_\theta},\, A \sim \pi_\theta(\cdot \mid S)}\!\left[\nabla_\theta \log \pi_\theta(A \mid S)\, Q^{\pi_\theta}(S,A)\right].
\end{equation}
\end{theorem}

\begin{proof}
Начинаем с~\eqref{eq:rtg-pg}. По линейности математического ожидания,
\[
\nabla_\theta J(\theta)
= \sum_{t=0}^{T-1} \gamma^t\, \E\!\left[\nabla_\theta \log \pi_\theta(A_t \mid S_t)\, G_t\right].
\]
Применяем теорему о полном математическом ожидании, обусловливаясь на $(S_t, A_t)$:
\begin{align*}
&\E\!\left[\nabla_\theta \log \pi_\theta(A_t \mid S_t)\, G_t\right]\\
&\quad= \E\!\left[\nabla_\theta \log \pi_\theta(A_t \mid S_t)\, \E[G_t \mid S_t, A_t]\right]\\
&\quad= \E\!\left[\nabla_\theta \log \pi_\theta(A_t \mid S_t)\, Q^{\pi_\theta}(S_t, A_t)\right].
\end{align*}
Раскрывая внешнее математическое ожидание как сумму по $(s,a)$ и используя тождество $\pi_\theta(a \mid s)\, \nabla_\theta \log \pi_\theta(a \mid s) = \nabla_\theta \pi_\theta(a \mid s)$:
\[
\nabla_\theta J(\theta)
= \sum_{t=0}^{T-1} \gamma^t \sum_s \Prob(S_t = s) \sum_a \nabla_\theta \pi_\theta(a \mid s)\, Q^{\pi_\theta}(s,a).
\]
Собирая коэффициент при каждом состоянии $s$, получаем $\sum_t \gamma^t \Prob(S_t = s) = d^{\gamma}_{\pi_\theta}(s)$, что соответствует определению~\ref{def:visitation}, откуда следует результат.
\end{proof}

\begin{remark}
\label{rem:pgt-model-free}
Теорема~\ref{thm:pgt} примечательна: градиент $J(\theta)$ по параметрам стратегии зависит от динамики переходов \emph{только через} функцию ценности действий $Q^{\pi_\theta}$\,---\,и даже её можно оценить по выборкам, не зная $p$. Именно это делает методы градиента стратегии истинно свободными от модели.
\end{remark}

% ===========================================================
\section{REINFORCE}
\label{sec:reinforce}

\index{REINFORCE}

\subsection{Алгоритм}

REINFORCE (Уильямс, 1992) \,---\, прямая реализация метода Монте-Карло для теоремы~\ref{thm:rtg-pg}. Поскольку $\E_{\pi_\theta}[G_t \mid S_t, A_t] = Q^{\pi_\theta}(S_t, A_t)$, мы можем подставить выборочную отдачу $G_t$ вместо истинного $Q^{\pi_\theta}$:
\begin{equation}
\label{eq:reinforce-update}
\theta \leftarrow \theta + \alpha \sum_{t=0}^{T-1} \gamma^t G_t\, \nabla_\theta \log \pi_\theta(A_t \mid S_t).
\end{equation}
Это обновление является несмещённым шагом стохастического градиентного подъёма по $J(\theta)$.

\begin{algorithm}[ht]
\caption{REINFORCE (Уильямс, 1992)}
\label{alg:reinforce}
\KwIn{Стратегия $\pi_\theta$, размер шага $\alpha$, дисконт $\gamma$}
Инициализировать $\theta$ произвольно\;
\For{каждый эпизод}{
  Сгенерировать $\tau = (S_0, A_0, R_1, \ldots, S_{T-1}, A_{T-1}, R_T)$ следуя $\pi_\theta$\;
  \For{$t = 0, 1, \ldots, T-1$}{
    Вычислить $G_t \leftarrow \sum_{k=t}^{T-1} \gamma^{k-t} R_{k+1}$\;
    $\theta \leftarrow \theta + \alpha\, \gamma^t G_t\, \nabla_\theta \log \pi_\theta(A_t \mid S_t)$\;
  }
}
\end{algorithm}

Стоит сделать несколько замечаний по реализации. Во-первых, множитель $\gamma^t$ перед обновлением уменьшает размер шага для более поздних шагов по времени, отражая более низкий дисконтированный вес соответствующего члена градиента. Некоторые реализации опускают этот множитель; это соответствует оптимизации несколько иной (недисконтированной) целевой функции, но является распространённой практикой. Во-вторых, поскольку все обновления в эпизоде могут быть вычислены в конце эпизода, цикл по $t$ можно векторизовать: вычислить все значения $G_t$ за один обратный проход, затем просуммировать взвешенные градиенты. В-третьих, в случае стратегии softmax по действиям, $\nabla_\theta \log \pi_\theta(a \mid s)$ является функцией счёта, равной $\phi(s,a) - \sum_{a'} \pi_\theta(a' \mid s)\, \phi(s,a')$ для векторов признаков $\phi(s,a)$.

\subsection{Разобранный пример}

\begin{example}[REINFORCE на двухсостоянном МПР]
\label{ex:reinforce-two-state}
\index{REINFORCE!worked example}

Рассмотрим следующий простой эпизодический МПР:
\begin{itemize}[leftmargin=*]
\item Два состояния: $\cS = \{s_1, s_2\}$; два действия: $\cA = \{a_L, a_R\}$.
\item Начальное состояние всегда $s_1$. Длина эпизода $T = 2$.
\item Переходы: из $s_1$ действие $a_L$ ведёт в $s_2$ с вознаграждением $R=0$; действие $a_R$ остаётся в $s_1$ с вознаграждением $R=0$. Из $s_2$ (или $s_1$ при $t=1$) получается фиксированное терминальное вознаграждение $+1$ за $a_R$ и $0$ за $a_L$.
\item Стратегия: $\pi_\theta(a_R \mid s_1) = \sigma(\theta_1)$, $\pi_\theta(a_R \mid s_2) = \sigma(\theta_2)$, где $\sigma$ \,---\, сигмоидная функция.
\end{itemize}

Задать $\gamma = 1$, $\alpha = 0.1$ и инициализировать $\theta_1 = \theta_2 = 0$ (равномерная стратегия, вероятность каждого действия $0.5$).

\medskip
\noindent\textbf{Эпизод:}  Предположим, что выборочная траектория равна $(s_1, a_L, 0, s_2, a_R, 1)$.

\medskip
Вычислить отдачу с текущего момента:
\[
G_0 = 0 + 1 = 1, \qquad G_1 = 1.
\]

Градиенты функции счёта (при $\theta_1 = \theta_2 = 0$, т.\,е. $\pi = 0.5$ для обоих действий):
\[
\nabla_{\theta_1} \log \pi_\theta(a_L \mid s_1)
= \frac{\partial}{\partial \theta_1} \log(1 - \sigma(\theta_1))
= -\sigma(\theta_1) = -0.5,
\]
\[
\nabla_{\theta_2} \log \pi_\theta(a_R \mid s_2)
= \sigma(\theta_2)(1 - \sigma(\theta_2)) \cdot \frac{1}{\sigma(\theta_2)} = 1 - 0.5 = 0.5.
\]

Обновления градиента ($\gamma = 1$):
\[
\Delta \theta_1 = \alpha \cdot \gamma^0 \cdot G_0 \cdot (-0.5) = 0.1 \cdot 1 \cdot 1 \cdot (-0.5) = -0.05,
\]
\[
\Delta \theta_2 = \alpha \cdot \gamma^1 \cdot G_1 \cdot 0.5 = 0.1 \cdot 1 \cdot 1 \cdot 0.5 = +0.05.
\]

После обновления: $\theta_1 = -0.05$ (делая $a_L$ чуть менее вероятным из $s_1$) и $\theta_2 = +0.05$ (делая $a_R$ более вероятным из $s_2$). Интуитивно алгоритм работает правильно: траектория, давшая положительную отдачу, подталкивает стратегию к действиям, которые её породили.
\end{example}

\begin{remark}
REINFORCE сходится к локальному максимуму $J(\theta)$ при стандартных условиях стохастической аппроксимации (убывающие размеры шагов, $\sum \alpha_t = \infty$, $\sum \alpha_t^2 < \infty$). Однако его дисперсия может быть очень большой, поскольку $G_t$ может сильно варьироваться от эпизода к эпизоду. Следующие два раздела рассматривают снижение дисперсии.
\end{remark}

% ===========================================================
\section{Снижение дисперсии: базовые линии}
\label{sec:baselines}

\index{baseline}

\subsection{Теорема о базовой линии}

\textbf{Базовая линия} $b(s)$ \,---\, это любая функция состояния (но не действия), которая вычитается из отдачи в оценщике градиента стратегии. Ключевой результат состоит в том, что это не вносит смещения.

\begin{theorem}[Теорема о базовой линии]
\label{thm:baseline}
\index{baseline theorem}
Для любой функции $b : \cS \to \R$,
\begin{equation}
\label{eq:baseline-theorem}
\nabla_\theta J(\theta)
= \E_{\tau \sim p_\theta}\!\left[\sum_{t=0}^{T-1} \gamma^t \nabla_\theta \log \pi_\theta(A_t \mid S_t)\, \bigl(G_t - b(S_t)\bigr)\right].
\end{equation}
\end{theorem}

\begin{proof}
Достаточно показать, что дополнительный член, введённый $b$, имеет нулевое математическое ожидание. Для каждого фиксированного $t$,
\begin{align*}
\E\!\left[\nabla_\theta \log \pi_\theta(A_t \mid S_t)\, b(S_t)\right]
&= \E\!\left[b(S_t)\, \E\!\left[\nabla_\theta \log \pi_\theta(A_t \mid S_t) \mid S_t\right]\right] \\
&= \E\!\left[b(S_t) \cdot 0\right] = 0,
\end{align*}
где внутреннее математическое ожидание обращается в нуль по тождеству функции счёта (см. доказательство леммы~\ref{lem:causality}). Суммируя по $t$ и применяя линейность математического ожидания, получаем~\eqref{eq:baseline-theorem}.
\end{proof}

\subsection{Оптимальная базовая линия}

Поскольку базовая линия не меняет среднее значение оценщика, мы вольны выбирать её так, чтобы \emph{минимизировать дисперсию}. Рассмотрим дисперсию одношагового члена $g_t(\theta) = \gamma^t \nabla_\theta \log \pi_\theta(A_t \mid S_t)\, (G_t - b)$ (зависимость от состояния опущена для ясности). Дисперсия скалярной случайной величины $X$ удовлетворяет $\Var[X - c] = \Var[X] - 2c\, \Cov[X, 1] + c^2\, \Var[1]$; для вектора $\mathbf{v} = \mathbf{u}\,(G_t - b)$ аналогичное вычисление, применённое к каждой компоненте, даёт оптимальную базовую линию для $k$-й компоненты градиента:
\begin{equation}
\label{eq:optimal-baseline}
b^*_k(s) = \frac{\E\!\left[(\nabla_\theta \log \pi_\theta(A_t \mid S_t))^2_k\, G_t \;\middle|\; S_t = s\right]}{\E\!\left[(\nabla_\theta \log \pi_\theta(A_t \mid S_t))^2_k \;\middle|\; S_t = s\right]}.
\end{equation}
Это отношение математических ожиданий, зависящее от индекса компоненты $k$. Оптимальная базовая линия, таким образом, является вектором и, как правило, вычислительно неприемлема. На практике вместо неё используется функция ценности состояния $V^{\pi_\theta}(s)$.

\subsection{Базовая линия на основе функции ценности состояния}

Наиболее естественной и широко используемой базовой линией является $b(S_t) = V^{\pi_\theta}(S_t)$. Вычитание функции ценности из отдачи даёт \textbf{функцию преимущества}\index{advantage function}:
\begin{equation}
\label{eq:advantage-ch8}
A^{\pi_\theta}(s,a) = Q^{\pi_\theta}(s,a) - V^{\pi_\theta}(s).
\end{equation}
Теорема о градиенте стратегии с базовой линией на основе функции ценности принимает вид
\begin{equation}
\label{eq:pg-advantage}
\nabla_\theta J(\theta)
= \E\!\left[\sum_{t=0}^{T-1} \gamma^t\, \nabla_\theta \log \pi_\theta(A_t \mid S_t)\, A^{\pi_\theta}(S_t, A_t)\right].
\end{equation}
Поскольку $\E_\pi[A^{\pi}(s,a) \mid s] = 0$ по определению, преимущество имеет нулевое среднее под стратегией, что часто резко снижает дисперсию оценщика. На практике $V^{\pi_\theta}$ неизвестно и само должно быть оценено; следующий раздел рассматривает это.

\begin{example}[Влияние базовой линии на дисперсию]
\label{ex:baseline-variance}
Рассмотрим бандитную (однострановую) среду с двумя действиями. Действие $a_1$ всегда даёт вознаграждение $10$; действие $a_2$ всегда даёт вознаграждение $10.1$. Без базовой линии сигнал градиента пропорционален $10$ или $10.1$\,---\,оба большие и почти равные. Оценщик с трудом различает чуть лучшее действие $a_2$ среди шума. При базовой линии $b = 10.05$ (среднее вознаграждение) сигнал градиента пропорционален $-0.05$ или $+0.05$, который центрирован вокруг нуля и имеет значительно меньшую абсолютную величину, существенно снижая дисперсию.
\end{example}

% ===========================================================
\section{REINFORCE с обученной базовой линией}
\label{sec:reinforce-baseline}

\index{REINFORCE!with baseline}

В предыдущем разделе базовая линия $b(s) = V^{\pi_\theta}(s)$ считалась известной. На практике её необходимо оценивать. Простейший подход\,---\,поддерживать отдельную параметризованную функцию ценности $\hat V_w(s)$ (\emph{критика}) рядом со стратегией (\emph{актором}) и обучать обоих одновременно.

\begin{algorithm}[ht]
\caption{REINFORCE с обученной базовой линией}
\label{alg:reinforce-baseline}
\KwIn{Стратегия $\pi_\theta$, функция ценности $\hat V_w$, размеры шагов $\alpha_\theta, \alpha_w$, дисконт $\gamma$}
Инициализировать $\theta, w$ произвольно\;
\For{каждый эпизод}{
  Сгенерировать $\tau = (S_0, A_0, R_1, \ldots, S_{T-1}, A_{T-1}, R_T)$ следуя $\pi_\theta$\;
  \For{$t = 0, 1, \ldots, T-1$}{
    $G_t \leftarrow \sum_{k=t}^{T-1} \gamma^{k-t} R_{k+1}$\;
    $\delta_t \leftarrow G_t - \hat V_w(S_t)$ \tcp*{Оценка преимущества}
    $w \leftarrow w + \alpha_w\, \delta_t\, \nabla_w \hat V_w(S_t)$ \tcp*{Обновление функции ценности}
    $\theta \leftarrow \theta + \alpha_\theta\, \gamma^t\, \delta_t\, \nabla_\theta \log \pi_\theta(A_t \mid S_t)$ \tcp*{Обновление стратегии}
  }
}
\end{algorithm}

Обновление функции ценности минимизирует среднеквадратичную ошибку $\E[(G_t - \hat V_w(S_t))^2]$, что является задачей обучения с учителем. Заметим, что градиент по $w$ равен $\nabla_w \hat V_w(S_t)$, а целью служит выборочная отдача $G_t$ (то есть без бутстрэппинга), поэтому обновление функции ценности является обычным шагом градиента\,---\,а не полуградиентным шагом, как в TD(0) (см. главу~\ref{ch:td}).

\begin{remark}
Введение отдельного критика открывает путь к полному семейству актор-критик, где критик использует TD-бутстрэппинг вместо отдачи Монте-Карло. Это обменивает несмещённость REINFORCE на сниженную дисперсию ценой введения смещения из-за ошибки аппроксимации функции. Методы актор-критик рассматриваются в главе~\ref{ch:actor-critic}.
\end{remark}

Распространённое опасение состоит в том, не вносит ли обученная базовая линия смещение. Нет: ключевое состоит в том, что $\hat V_w(S_t)$ зависит только от состояния $S_t$, а не от действия $A_t$, выбранного на этом шаге. Поэтому теорема~\ref{thm:baseline} применима независимо от того, как получено $\hat V_w$. Оценка градиента стратегии $\delta_t\, \nabla_\theta \log \pi_\theta(A_t \mid S_t)$ остаётся несмещённой для $\nabla_\theta J(\theta)$ даже при $\hat V_w \neq V^{\pi_\theta}$, хотя неоптимальное $\hat V_w$ не будет снижать дисперсию так же эффективно, как истинная функция ценности.

% ===========================================================
\section{RLOO: REINFORCE Leave-One-Out}
\label{sec:rloo}

\index{RLOO}
\index{REINFORCE Leave-One-Out|see{RLOO}}

\subsection{Мотивация: градиент стратегии в контексте языковых моделей}

Современные пайплайны выравнивания больших языковых моделей (LLM) (см. главу~\ref{ch:practical-rlhf}) дообучают языковую модель $\pi_\theta$ с использованием сигнала вознаграждения $r(\mathbf{y} \mid \mathbf{x})$, который оценивает сгенерированное дополнение $\mathbf{y}$ к промпту $\mathbf{x}$. В этом контексте <<действие>> \,---\, это целое дополнение (или, эквивалентно, последовательность действий на уровне токенов), а <<состояние>> \,---\, промпт. Таким образом, возникает задача градиента стратегии с:
\begin{itemize}[leftmargin=*]
\item Очень высокой дисперсией на выборку (дополнения могут сильно различаться).
\item Отсутствием доступа к отдельно обученному критику (обучение головы функции ценности добавляет потребление памяти и нестабильность).
\item Естественной пакетной структурой: для каждого промпта $\mathbf{x}$ в мини-пакете модель генерирует $K$ независимых дополнений $\mathbf{y}^{(1)}, \ldots, \mathbf{y}^{(K)} \sim \pi_\theta(\cdot \mid \mathbf{x})$.
\end{itemize}
Ключом является пакетная структура. RLOO использует её для построения базовой линии с низкой дисперсией, несмещённой и не требующей критика.

\subsection{Оценщик RLOO}

\index{RLOO!estimator}

Зафиксируем промпт $\mathbf{x}$ и выберем $K$ н.о.р.\ дополнений $\mathbf{y}^{(1)}, \ldots, \mathbf{y}^{(K)} \sim \pi_\theta(\cdot \mid \mathbf{x})$ с соответствующими вознаграждениями $r_i = r(\mathbf{y}^{(i)} \mid \mathbf{x})$. \textbf{Базовая линия RLOO} для выборки $i$ \,---\, среднее вознаграждение остальных $K-1$ выборок:
\begin{equation}
\label{eq:rloo-baseline}
b_i = \frac{1}{K-1} \sum_{j \neq i} r_j = \frac{1}{K-1}\left(\sum_{j=1}^K r_j - r_i\right).
\end{equation}
\textbf{Оценщик градиента стратегии RLOO} тогда равен
\begin{equation}
\label{eq:rloo-estimator}
\hat g_{\text{RLOO}} = \frac{1}{K} \sum_{i=1}^K (r_i - b_i)\, \nabla_\theta \log \pi_\theta(\mathbf{y}^{(i)} \mid \mathbf{x}),
\end{equation}
где $\nabla_\theta \log \pi_\theta(\mathbf{y}^{(i)} \mid \mathbf{x}) = \sum_{t} \nabla_\theta \log \pi_\theta(y^{(i)}_t \mid \mathbf{x}, y^{(i)}_{<t})$ \,---\, суммарная функция счёта по токенам.

Подставляя~\eqref{eq:rloo-baseline} в~\eqref{eq:rloo-estimator}, можно также записать
\begin{equation}
\label{eq:rloo-expanded}
\hat g_{\text{RLOO}} = \frac{1}{K} \sum_{i=1}^K \left(r_i - \frac{\bar r - r_i / K}{(K-1)/K}\right) \nabla_\theta \log \pi_\theta(\mathbf{y}^{(i)} \mid \mathbf{x}),
\end{equation}
где $\bar r = \frac{1}{K}\sum_j r_j$. Более компактное выражение:
\begin{equation}
\label{eq:rloo-clean}
r_i - b_i = r_i - \frac{\sum_{j \neq i} r_j}{K-1} = \frac{K\, r_i - \sum_{j=1}^K r_j}{K-1} = \frac{K}{K-1}\!\left(r_i - \bar r\right).
\end{equation}
Таким образом, преимущество для выборки $i$ является масштабированным отклонением $r_i$ от среднего по пакету, с поправочным коэффициентом $K/(K-1)$, учитывающим структуру leave-one-out.

\begin{algorithm}[ht]
\caption{Обновление градиента стратегии RLOO}
\label{alg:rloo}
\KwIn{Стратегия $\pi_\theta$, функция вознаграждения $r$, размер пакета $K$, размер шага $\alpha$}
\For{каждый промпт $\mathbf{x}$ в мини-пакете}{
  Выбрать $K$ дополнений $\mathbf{y}^{(1)}, \ldots, \mathbf{y}^{(K)} \sim \pi_\theta(\cdot \mid \mathbf{x})$\;
  Вычислить вознаграждения $r_i \leftarrow r(\mathbf{y}^{(i)} \mid \mathbf{x})$ для $i = 1, \ldots, K$\;
  Вычислить среднее по пакету $\bar r \leftarrow \frac{1}{K}\sum_{i=1}^K r_i$\;
  \For{$i = 1, \ldots, K$}{
    $b_i \leftarrow \frac{K}{K-1}(\bar r - r_i/K) \cdot \frac{1}{1} = \frac{1}{K-1}\bigl(\sum_j r_j - r_i\bigr)$ \tcp*{Базовая линия LOO}
    $\hat A_i \leftarrow r_i - b_i = \frac{K}{K-1}(r_i - \bar r)$\;
  }
  $\theta \leftarrow \theta + \alpha \cdot \frac{1}{K}\sum_{i=1}^K \hat A_i\, \nabla_\theta \log \pi_\theta(\mathbf{y}^{(i)} \mid \mathbf{x})$\;
}
\end{algorithm}

\subsection{Несмещённость RLOO}

\begin{theorem}[RLOO несмещён]
\label{thm:rloo-unbiased}
\index{RLOO!unbiasedness}
Оценщик RLOO $\hat g_{\text{RLOO}}$, определённый в~\eqref{eq:rloo-estimator}, является несмещённым оценщиком $\nabla_\theta J(\theta)$:
\[
\E\!\left[\hat g_{\text{RLOO}}\right] = \nabla_\theta J(\theta).
\]
\end{theorem}

\begin{proof}
Достаточно показать, что $b_i$ удовлетворяет условиям теоремы о базовой линии (теорема~\ref{thm:baseline}), т.\,е. что она не зависит от действия $\mathbf{y}^{(i)}$ при данном состоянии (промпте $\mathbf{x}$). Базовая линия $b_i = \frac{1}{K-1}\sum_{j \neq i} r_j$ является функцией $K-1$ остальных дополнений $\mathbf{y}^{(j)}$, $j \neq i$. Поскольку все $K$ дополнений выбираются независимо из $\pi_\theta(\cdot \mid \mathbf{x})$, базовая линия $b_i$ не зависит от $\mathbf{y}^{(i)}$ условно на $\mathbf{x}$.

Формально, взяв математическое ожидание $i$-го члена:
\begin{align*}
\E\!\left[(r_i - b_i)\, \nabla_\theta \log \pi_\theta(\mathbf{y}^{(i)} \mid \mathbf{x})\right]
&= \E\!\left[r_i\, \nabla_\theta \log \pi_\theta(\mathbf{y}^{(i)} \mid \mathbf{x})\right]
 - \E\!\left[b_i\, \nabla_\theta \log \pi_\theta(\mathbf{y}^{(i)} \mid \mathbf{x})\right].
\end{align*}
Для второго члена обусловимся на все $\mathbf{y}^{(j)}$ с $j \neq i$ (что определяет $b_i$):
\begin{align*}
\E\!\left[b_i\, \nabla_\theta \log \pi_\theta(\mathbf{y}^{(i)} \mid \mathbf{x})\right]
&= \E\!\left[b_i\, \E\!\left[\nabla_\theta \log \pi_\theta(\mathbf{y}^{(i)} \mid \mathbf{x}) \;\middle|\; \mathbf{x},\, \{r_j\}_{j \neq i}\right]\right].
\end{align*}
Внутреннее математическое ожидание \,---\, это в точности тождество функции счёта, применённое к распределению $\pi_\theta(\cdot \mid \mathbf{x})$:
\[
\E_{\mathbf{y}^{(i)} \sim \pi_\theta(\cdot \mid \mathbf{x})}\!\left[\nabla_\theta \log \pi_\theta(\mathbf{y}^{(i)} \mid \mathbf{x})\right] = \nabla_\theta \sum_{\mathbf{y}} \pi_\theta(\mathbf{y} \mid \mathbf{x}) = \nabla_\theta 1 = 0.
\]
Поэтому вычтенный член с базовой линией равен нулю, и
\[
\E\!\left[(r_i - b_i)\, \nabla_\theta \log \pi_\theta(\mathbf{y}^{(i)} \mid \mathbf{x})\right]
= \E\!\left[r_i\, \nabla_\theta \log \pi_\theta(\mathbf{y}^{(i)} \mid \mathbf{x})\right]
= \nabla_\theta J(\theta)\big|_{\mbox{одна выборка}}.
\]
Усредняя по $i = 1, \ldots, K$, получаем $\hat g_{\text{RLOO}}$, математическое ожидание которого равно $\nabla_\theta J(\theta)$.
\end{proof}

\subsection{Снижение дисперсии в RLOO}

\index{RLOO!variance}

Чтобы понять, почему RLOO снижает дисперсию по сравнению с обычным REINFORCE, сравним два оценщика. Обычный REINFORCE использует полную отдачу $r_i$ в качестве веса градиента (без базовой линии):
\[
\hat g_{\text{RF}} = \frac{1}{K}\sum_{i=1}^K r_i\, \nabla_\theta \log \pi_\theta(\mathbf{y}^{(i)} \mid \mathbf{x}).
\]
Пусть $g_i = \nabla_\theta \log \pi_\theta(\mathbf{y}^{(i)} \mid \mathbf{x})$ и $\mu = \E[r_i]$. Тогда $\Var[r_i\, g_i] \approx \Var[r_i] \Var[g_i] + \mu^2 \Var[g_i]$ (без учёта перекрёстных членов), что растёт с величиной $\mu$. RLOO заменяет $r_i$ на $r_i - b_i \approx r_i - \mu$ (с точностью до поправки $K/(K-1)$), среднее значение которого равно нулю. Доминирующий член $\mu^2$ тем самым устраняется, и дисперсия масштабируется с $\Var[r_i]$, а не с $\E[r_i^2]$.

Точнее, при большом $K$ базовая линия RLOO $b_i \to \mu = \E_{\pi_\theta}[r]$ (по закону больших чисел), и RLOO приближается к оракульной базовой линии $b = \mu$. Остаточная дисперсия $r_i - b_i$ равна $\Var[r_i] / (K-1)$, по сравнению с $\Var[r_i] + \mu^2$ при отсутствии базовой линии. Когда вознаграждения имеют большое положительное среднее ($\mu \gg 0$), снижение дисперсии может быть значительным.

\subsection{Связь с RLHF}

\index{RLOO!RLHF connection}
\index{RLHF!RLOO}

В пайплайнах RLHF (глава~\ref{ch:practical-rlhf}) модель вознаграждения $r_\phi(\mathbf{y} \mid \mathbf{x})$ обеспечивает скалярную обратную связь по дополнениям модели. RLOO хорошо вписывается: для каждого промпта $\mathbf{x}$ выбирается $K$ дополнений из текущей стратегии $\pi_\theta$, они оцениваются моделью вознаграждения, и вычисляется градиент RLOO. Сеть-критик не нужна, что экономит память GPU и устраняет нестабильность, которая может возникать при совместном обучении критика и стратегии.

RLOO был введён Ахмадяном и~др.\ (2024) и показал превосходство над базовыми методами RLHF на основе PPO при одинаковом количестве запросов к модели вознаграждения, в основном благодаря снижению дисперсии от базовой линии leave-one-out и устранению критика. Его также можно комбинировать со штрафом за KL-дивергенцию от референсной стратегии $\pi_{\text{ref}}$, заменив $r_i$ на $r_i - \beta \log(\pi_\theta(\mathbf{y}^{(i)}) / \pi_{\text{ref}}(\mathbf{y}^{(i)}))$, что контролирует, насколько сильно дообученная модель отклоняется от предобученной.

\begin{remark}
\label{rem:rloo-grpo}
RLOO тесно связан с Group Relative Policy Optimisation (GRPO), который также использует пакет выборок на промпт и нормализует преимущества внутри группы. Ключевое отличие в том, что GRPO делит на стандартное отклонение внутри группы (стандартизируя преимущество), тогда как RLOO использует немасштабированное отклонение leave-one-out. Оба принадлежат к семейству оценщиков с <<самоопорными базовыми линиями>>, которые извлекают снижающий дисперсию сигнал из самого пакета, без отдельно обученного критика.
\end{remark}

% ===========================================================
\section{Натуральный градиент стратегии}
\label{sec:npg}

\index{natural policy gradient}

\subsection{Проблема обычного градиента}

Обычный градиент $\nabla_\theta J(\theta)$ определён относительно евклидовой метрики в пространстве параметров. Шаг градиента $\theta \leftarrow \theta + \alpha \nabla_\theta J(\theta)$ перемещает на фиксированное евклидово расстояние в пространстве $\theta$, но соответствующее изменение распределения стратегии $\pi_\theta$ зависит от локальной геометрии многообразия стратегий. Вблизи детерминированной стратегии (где вероятность некоторого действия близка к 1) малый шаг в $\theta$ вызывает большое изменение $\pi_\theta$; вблизи равномерной стратегии тот же шаг вызывает меньшее распределительное изменение. Обычный градиент поэтому плохо откалиброван: его величина отражает евклидово расстояние в пространстве параметров, а не изменение распределения.

\subsection{Матрица информации Фишера}

\index{Fisher information matrix}

Натуральный градиент стратегии исправляет это, используя \textbf{матрицу информации Фишера}\index{Fisher information matrix} стратегии в качестве римановой метрики:
\begin{equation}
\label{eq:fisher}
F(\theta) = \E_{S \sim d^{\gamma}_{\pi_\theta},\, A \sim \pi_\theta(\cdot \mid S)}\!\left[\nabla_\theta \log \pi_\theta(A \mid S)\, \nabla_\theta \log \pi_\theta(A \mid S)^\top\right].
\end{equation}
Эта матрица измеряет ожидаемое внешнее произведение векторов функции счёта; она фиксирует, насколько чувствительно распределение $\pi_\theta$ реагирует на малые изменения $\theta$. Эквивалентно, $F(\theta)$ является гессианом KL-дивергенции $\KL{\pi_\theta}{\pi_{\theta'}}$ по $\theta'$, вычисленным при $\theta' = \theta$.

\textbf{Натуральный градиент}\index{natural gradient} \,---\, это направление наискорейшего подъёма в пространстве распределений, измеренное KL-дивергенцией:
\begin{equation}
\label{eq:natural-gradient}
\tilde\nabla_\theta J(\theta) = F(\theta)^{-1} \nabla_\theta J(\theta).
\end{equation}
Обновление по натуральному градиенту стратегии:
\begin{equation}
\label{eq:npg-update}
\theta \leftarrow \theta + \alpha\, F(\theta)^{-1} \nabla_\theta J(\theta).
\end{equation}

\begin{theorem}[Амари, 1998; Какаде, 2002]
\label{thm:npg}
Натуральный градиент $\tilde\nabla_\theta J(\theta) = F(\theta)^{-1} \nabla_\theta J(\theta)$ является направлением наискорейшего подъёма для $J$ при ограничении, что KL-дивергенция между старой и новой стратегией удовлетворяет $\KL{\pi_\theta}{\pi_{\theta + \Delta\theta}} \leq \epsilon$, для малого $\epsilon$.
\end{theorem}

\begin{proof}[Набросок доказательства]
Задача условной оптимизации разложения Тейлора первого порядка $J(\theta + \Delta\theta) \approx J(\theta) + \nabla_\theta J(\theta)^\top \Delta\theta$ при ограничении $\Delta\theta^\top F(\theta) \Delta\theta \leq \epsilon$ (ограничение в форме шара KL в пространстве параметров) даёт по неравенству Коши-Шварца во взвешенном $F$ скалярном произведении оптимальное направление $\Delta\theta^* \propto F(\theta)^{-1} \nabla_\theta J(\theta)$.
\end{proof}

\subsection{Практические соображения}

Прямое вычисление $F(\theta)^{-1}$ требует $O(d^2)$ памяти и $O(d^3)$ времени, что неприемлемо для больших нейронных сетей. На практике используются несколько приближений:
\begin{itemize}[leftmargin=*]
\item \textbf{Метод сопряжённых градиентов.}  Решать $F(\theta)\, \tilde g = \nabla_\theta J(\theta)$ итерационно. Каждый шаг требует произведения Фишера на вектор, которое может быть вычислено за $O(d)$ с помощью обратного распространения через функцию счёта. Это подход, используемый в TRPO (Шульман и~др., 2015).

\item \textbf{Кронекерово-факторизованное приближение (K-FAC).}  Аппроксимировать $F(\theta)$ как блочно-диагональную матрицу с кронекеровой структурой, используя послойную факторизацию нейронных сетей (Мартенс и Гросс, 2015).

\item \textbf{Proximal Policy Optimization (PPO).}  Вместо явного обращения Фишера, PPO использует ограниченную суррогатную целевую функцию для неявного ограничения KL-дивергенции, достигая аналогичного поведения при значительно меньших затратах (Шульман и~др., 2017).
\end{itemize}
Натуральный градиент стратегии концептуально важен, поскольку обеспечивает теоретическую основу для методов доверительной области и проксимальных методов; его практические реализации являются рабочими лошадками современного глубокого ОП и дообучения LLM.

% ===========================================================
\section{Анализ сходимости}
\label{sec:pg-convergence}

\index{REINFORCE!convergence}

\subsection{Сходимость REINFORCE}

REINFORCE \,---\, это алгоритм стохастического градиентного подъёма, применённый к $J(\theta)$. Его свойства сходимости следуют из общей теории стохастической аппроксимации (Роббинс и Монро, 1951), при условии, что оценщик градиента несмещён (установлено в теореме~\ref{thm:rtg-pg}) и размеры шагов удовлетворяют условиям Роббинса-Монро:
\begin{equation}
\label{eq:rm-conditions-ch8}
\sum_{k=1}^\infty \alpha_k = \infty, \qquad \sum_{k=1}^\infty \alpha_k^2 < \infty.
\end{equation}

\begin{theorem}[Сходимость REINFORCE]
\label{thm:reinforce-convergence}
Предположим, что $\pi_\theta(a \mid s)$ дважды непрерывно дифференцируема по $\theta$, и что градиент $\nabla_\theta \log \pi_\theta$ ограничен. Если размеры шагов удовлетворяют~\eqref{eq:rm-conditions-ch8} и $J(\theta)$ ограничена сверху, то итерации REINFORCE удовлетворяют
\[
\liminf_{k \to \infty} \|\nabla_\theta J(\theta_k)\|^2 = 0 \quad \mbox{п.н.}
\]
То есть REINFORCE сходится к стационарной точке $J$.
\end{theorem}

Теорема гарантирует сходимость к \emph{стационарной точке}, которая может быть локальным максимумом, седловой точкой или (в вырожденных случаях) локальным минимумом. Глобальная оптимальность не гарантируется для невыпуклого $J$. Однако для \emph{табличных} стратегий (без аппроксимации функции, конечные пространства состояний и действий) целевая функция градиента стратегии не имеет ложных локальных максимумов при мягких условиях (Агарвал и~др., 2021), поэтому сходимость к глобальному оптимуму может быть гарантирована.

\subsection{Скорость сходимости}

Для невыпуклого случая современное состояние теории \,---\, скорость $\tilde O(1/\sqrt{K})$ сходимости к стационарной точке по числу траекторий $K$:

\begin{theorem}[Скорость сходимости, неформально]
\label{thm:reinforce-rate}
При стандартных условиях регулярности REINFORCE с постоянным размером шага $\alpha$ находит точку $\theta$ с $\|\nabla_\theta J(\theta)\|^2 \leq \epsilon$ не более чем за $O(\sigma^2 / (\alpha \epsilon))$ эпизодов, где $\sigma^2$ \,---\, верхняя оценка дисперсии оценщика градиента.
\end{theorem}

Эта скорость раскрывает центральную роль дисперсии: снижение $\sigma^2$ (через базовые линии или RLOO) непосредственно снижает сложность по выборкам. При базовой линии, снижающей дисперсию в $c$ раз, требуется в $c$ раз меньше эпизодов для достижения той же точности.

\subsection{Компромисс между смещением и дисперсией}

Все оценщики градиента стратегии предполагают компромисс между смещением и дисперсией в том, как они оценивают $Q^{\pi_\theta}(s,a)$:
\begin{itemize}[leftmargin=*]
\item \textbf{Отдача Монте-Карло $G_t$:} несмещённая, но с высокой дисперсией (накапливает шум вознаграждения на протяжении всего эпизода).
\item \textbf{Цель TD(0) $R_{t+1} + \gamma \hat V(S_{t+1})$:} смещённая (из-за ошибки аппроксимации $\hat V$), но с низкой дисперсией (только один шаг шума).
\item \textbf{$n$-шаговая отдача:} интерполирует между двумя крайностями; увеличение $n$ снижает смещение и увеличивает дисперсию.
\item \textbf{GAE (обобщённая оценка преимущества, раздел~\ref{sec:gae}) с $\lambda \in [0,1]$:} экспоненциально взвешенное среднее $n$-шаговых отдач, сочетающее низкое смещение ($\lambda$ близко к 1) с низкой дисперсией ($\lambda$ близко к 0).
\end{itemize}

Оценщик REINFORCE соответствует крайнему случаю $n = T$ (полный Монте-Карло). RLOO сохраняет отдачу Монте-Карло, но снижает дисперсию через базовую линию leave-one-out, тогда как методы актор-критик (глава~\ref{ch:actor-critic}) снижают дисперсию через бутстрэппинг ценой некоторого смещения.

\begin{remark}
\label{rem:variance-dominates}
На практике дисперсия обычно является доминирующей проблемой для методов градиента стратегии, особенно в задачах с длинным горизонтом или с большими масштабами вознаграждений. Снижения базовой линии и отдачи с текущего момента вместе часто уменьшают дисперсию на один-два порядка величины по сравнению с наивным оценщиком~\eqref{eq:basic-pg}, что составляет разницу между сходимостью и расходимостью.
\end{remark}

% ===========================================================
\section{Сравнение оценщиков градиента стратегии}
\label{sec:pg-comparison}

\index{policy gradient!estimators comparison}

В таблице~\ref{tab:pg-comparison-ch8} обобщены свойства оценщиков, разработанных в данной главе. Используются следующие обозначения: $\hat g$ обозначает оценщик градиента, $b$ \,---\, базовую линию, $K$ \,---\, число выборок на промпт/состояние.

\begin{table}[ht]
\centering
\caption{Сравнение оценщиков градиента стратегии.}
\label{tab:pg-comparison-ch8}
\small
\begin{tabular}{@{}llllp{3.2cm}@{}}
\toprule
\textbf{Оценщик} & \textbf{Базовая линия} & \textbf{Смещ./Дисп.} & \textbf{Критик?} & \textbf{Примечания} \\
\midrule
REINFORCE & Нет & Несмещ./Очень выс. & Нет & Простейший; редко используется отдельно \\
\addlinespace
REINFORCE+$b$ & Фикс. $b$ & Несмещ./Высок. & Нет & Снижает среднее; неоптимален \\
\addlinespace
REINFORCE+$\hat V$ & $\hat V_w(s)$ & Несмещ./Средн. & Да & Наиболее распространённая б.л. \\
\addlinespace
RLOO & LOO-среднее & Несмещ./Низк. & Нет & $K\!\geq\!2$ выборки; без критика \\
\addlinespace
Актор-критик & $\hat V_w(s)$ & Смещ./Низк. & Да & Смещение из-за бутстрэппинга \\
\addlinespace
Натур. ГС & Нет & Несмещ./Средн. & Нет & Лучшая геометрия \\
\addlinespace
TRPO/PPO & $\hat V_w(s)$ & Смещ./Низк. & Да & KL/отсечение; лучший результат \\
\bottomrule
\end{tabular}
\end{table}

Из таблицы~\ref{tab:pg-comparison-ch8} стоит сделать несколько наблюдений.

Во-первых, все оценщики в столбце <<Смещение>>, отмеченные как <<Несмещённые>>, являются несмещёнными оценщиками $\nabla_\theta J(\theta)$: градиент стратегии является точным градиентом истинной целевой функции. Смещение возникает только при замене отдачи Монте-Карло на бутстрэппинговые цели (как в методах актор-критик) или когда ошибка аппроксимации функции ценности распространяется на оценку преимущества.

Во-вторых, столбец <<Критик?>> раскрывает фундаментальный выбор дизайна. Методы без критика (REINFORCE, RLOO) проще, но полагаются на отдачу Монте-Карло с высокой дисперсией; методы с критиком (актор-критик, PPO) достигают более низкой дисперсии ценой отдельной сети, дополнительных гиперпараметров (скорость обучения критика, количество эпох обновления) и потенциальной нестабильности при неточном критике.

В-третьих, RLOO достигает лучшего из обоих миров\,---\,без критика, но с существенно меньшей дисперсией, чем обычный REINFORCE,\,---\,используя естественную пакетную структуру обучения LLM. При одной выборке ($K=1$) RLOO сводится к обычному REINFORCE (базовая линия LOO при одной выборке равна нулю), поэтому улучшение зависит от наличия $K \geq 2$.

В-четвёртых, натуральный градиент стратегии и его практические потомки (TRPO, PPO) предпочтительны при необходимости больших обновлений параметров, поскольку они инвариантны к перепараметризации стратегии: направление обновления одинаково независимо от математического представления $\pi_\theta$.

% ===========================================================
\section*{Упражнения}
\addcontentsline{toc}{section}{Упражнения}

\begin{exercise}
\label{ex:pg-score}
Пусть $\pi_\theta(a \mid s) = \mathrm{softmax}(\theta_a)$ \,---\, табличная softmax-стратегия над $n$ действиями.
\begin{enumerate}
\item Покажите, что $\nabla_{\theta_a} \log \pi_\theta(a' \mid s) = \ind[a = a'] - \pi_\theta(a \mid s)$.
\item Проверьте тождество функции счёта: $\sum_{a} \pi_\theta(a \mid s)\, \nabla_\theta \log \pi_\theta(a \mid s) = 0$.
\item Пусть $n = 3$, $\theta = (1, 0, -1)$, и наблюдается действие $a = 2$ с отдачей $G = 5$. Вычислите обновление REINFORCE $\Delta\theta$ с размером шага $\alpha = 0.01$.
\end{enumerate}
\end{exercise}

\begin{exercise}
\label{ex:pg-causality}
Докажите лемму о причинности (лемма~\ref{lem:causality}) непосредственно из определения $p_\theta(\tau)$, не используя теорему о полном математическом ожидании. Подсказка: явно распишите сумму по траекториям и разложите её по шагам времени.
\end{exercise}

\begin{exercise}
\label{ex:pg-baseline-variance}
Рассмотрим однострановой двухдействийный бандит. Пусть $r(a_1) = 1$ и $r(a_2) = 3$ с вероятностью 1. Пусть $\pi_\theta(a_1) = \sigma(-\theta)$, $\pi_\theta(a_2) = \sigma(\theta)$.
\begin{enumerate}
\item Вычислите $J(\theta) = \E[r]$ как функцию от $\theta$.
\item Вычислите $\nabla_\theta J(\theta)$ аналитически.
\item Вычислите дисперсию оценщика градиента REINFORCE (без базовой линии) как функцию от $\theta$.
\item Покажите, что базовая линия $b = \E[r] = J(\theta)$ снижает дисперсию. Вычислите дисперсию с этой базовой линией.
\item Какая базовая линия является оптимальной согласно~\eqref{eq:optimal-baseline}?
\end{enumerate}
\end{exercise}

\begin{exercise}
\label{ex:rloo-unbiased-direct}
Докажите несмещённость оценщика RLOO непосредственно, не ссылаясь на теорему~\ref{thm:baseline}. Конкретно, для $K = 2$ дополнений $y^{(1)}, y^{(2)} \sim \pi_\theta(\cdot \mid x)$ с вознаграждениями $r_1, r_2$, покажите, что
\[
\E\!\left[(r_1 - r_2)\, \nabla_\theta \log \pi_\theta(y^{(1)} \mid x) + (r_2 - r_1)\, \nabla_\theta \log \pi_\theta(y^{(2)} \mid x)\right] = 2\, \nabla_\theta J(\theta).
\]
\end{exercise}

\begin{exercise}
\label{ex:rloo-variance}
В обозначениях раздела~\ref{sec:rloo} пусть $r_1, \ldots, r_K$ \,---\, н.о.р.\ со средним $\mu$ и дисперсией $\sigma^2$, а $g_1, \ldots, g_K$ \,---\, функции счёта (независимые от вознаграждений по тождеству функции счёта, поэтому в данном упражнении считайте их фиксированными). Покажите, что
\[
\Var\!\left[\frac{1}{K}\sum_{i=1}^K (r_i - b_i)\, g_i\right]
\approx \frac{\sigma^2}{K-1} \cdot \frac{1}{K}\sum_{i=1}^K g_i^2
\]
при большом $K$, и сравните это с дисперсией оценщика обычного REINFORCE $\frac{1}{K}\sum_i r_i g_i$. При каких условиях снижение дисперсии наибольшее?
\end{exercise}

\begin{exercise}
\label{ex:npg-fisher}
\begin{enumerate}
\item Пусть $\pi_\theta$ \,---\, гауссовская стратегия: $\pi_\theta(a \mid s) = \cN(a;\, \mu_\theta(s),\, \sigma^2)$ с фиксированной дисперсией $\sigma^2$. Вычислите матрицу информации Фишера $F(\theta)$ через $\nabla_\theta \mu_\theta(s)$.
\item Покажите, что когда $\mu_\theta(s) = \theta^\top \phi(s)$ линейно, натуральный градиент сводится к $\tilde\nabla_\theta J = \sigma^2\, F_\phi^{-1} \nabla_\theta J$, где $F_\phi$ \,---\, матрица ковариации признаков.
\item Объясните интуитивно, почему обновление по натуральному градиенту инвариантно к перепараметризации стратегии.
\end{enumerate}
\end{exercise}

\begin{exercise}
\label{ex:pg-continuous}
Рассмотрим МПР с непрерывным пространством действий $\cA = \R$ и гауссовской стратегией $\pi_\theta(a \mid s) = \cN(a;\, \mu_\theta(s),\, e^{2\rho})$, где $\theta = (\theta_\mu, \rho)$ параметризует сеть среднего и логарифмическое стандартное отклонение.
\begin{enumerate}
\item Выведите $\nabla_{\theta_\mu} \log \pi_\theta(a \mid s)$ и $\nabla_\rho \log \pi_\theta(a \mid s)$.
\item Интерпретируйте градиент по $\rho$: когда обновление увеличивает энтропию, а когда уменьшает?
\item Покажите, что обновление REINFORCE по $\rho$ реализует автоматическое регулирование энтропии исключительно через сигнал вознаграждения.
\end{enumerate}
\end{exercise}

\begin{exercise}
\label{ex:pg-convergence}
\begin{enumerate}
\item Сформулируйте условия Роббинса-Монро~\eqref{eq:rm-conditions-ch8} и приведите один пример последовательности размеров шагов, удовлетворяющей им, и один пример, не удовлетворяющей.
\item Рассмотрим REINFORCE с постоянным размером шага $\alpha > 0$. Почему это \emph{не} удовлетворяет условиям Роббинса-Монро, и к чему вместо этого сходится алгоритм?
\item Объясните компромисс между смещением и дисперсией между REINFORCE (отдача Монте-Карло) и методом актор-критик (цель TD) в контексте теоремы~\ref{thm:reinforce-rate}. Как снижение дисперсии через базовую линию влияет на оценку сложности по выборкам?
\end{enumerate}
\end{exercise}

\begin{exercise}
\label{ex:pg-lm}
В контексте языковой модели стратегия $\pi_\theta$ генерирует последовательности токенов авторегрессионно. Предположим, что мы используем RLOO с $K = 4$ дополнениями на промпт.
\begin{enumerate}
\item Выпишите преимущество RLOO $\hat A_i = r_i - b_i$ для $i = 1, 2, 3, 4$ по формуле из~\eqref{eq:rloo-clean}.
\item Предположим, что четыре вознаграждения равны $r = (0.2, 0.8, 0.5, 0.3)$. Вычислите все четыре преимущества.
\item Покажите, что $\sum_{i=1}^4 \hat A_i = 0$. Всегда ли это верно? Докажите.
\item Обсудите, почему свойство нулевой суммы преимуществ RLOO (при суммировании по пакету) может быть выгодным или невыгодным на практике.
\end{enumerate}
\end{exercise}

\begin{exercise}
\label{ex:pg-comparison-design}
Вы разрабатываете систему ОП для дообучения предобученной языковой модели на задаче генерации текста. Вам доступна функция вознаграждения типа <<чёрный ящик>> и возможность генерировать до $K = 8$ дополнений на промпт параллельно.
\begin{enumerate}
\item Сравните REINFORCE с обученной базовой линией функции ценности, RLOO и PPO по следующим измерениям: (a) требования к памяти, (b) количество гиперпараметров, (c) потенциальное смещение в оценке градиента, (d) простота реализации.
\item При каких обстоятельствах вы предпочли бы RLOO вместо PPO? При каких предпочли бы PPO?
\item Предложите простой эксперимент для эмпирического определения оптимального $K$ для RLOO в вашем контексте.
\end{enumerate}
\end{exercise}
